
This work intersects three research areas: RLHF methods that use preference data for training, alignment evaluation approaches, and embedding-based analysis techniques.

\subsubsection{2.1 RLHF and Preference Learning}

Christiano et al. [2017] introduced deep reinforcement learning from human preferences, demonstrating that pairwise comparisons provide scalable supervision. Ziegler et al. [2019] applied this framework to language model fine-tuning. These foundational works established the paradigm of learning reward functions from human feedback, treating preference data as training input.

Ouyang et al. [2022] scaled RLHF to InstructGPT with quality-controlled human annotations. Bai et al. [2022] released the HH-RLHF dataset with 160,000+ preference pairs and explicit helpfulness/harmlessness criteria. Both works train Bradley-Terry reward models on preference data, then discard the data after model training.

This work differs by testing whether the preference data itself encodes geometric structure exploitable for evaluation. The negative finding (no clustering in standard embeddings) establishes that preference data structure is implicit in learned reward representations rather than manifest in general-purpose semantic embeddings.

\subsubsection{2.2 Alignment Evaluation Methods}

Hendrycks et al. [2020] introduced the ETHICS benchmark with validated human annotations across multiple moral scenarios. Lin et al. [2022] developed TruthfulQA for evaluating truthfulness through curated question-answer pairs. These approaches require manual curation for each evaluation dimension.

This work tests whether existing RLHF preference data could serve as reusable benchmarks without additional annotation. The negative result shows this requires safety-specialized representations—pretrained embeddings are insufficient. The base-rate validation protocol (45.6\% genuine violations) demonstrates that RLHF datasets contain quality ground truth, though not in forms detectable by standard embeddings.

\subsubsection{2.3 Embedding-Based Safety Analysis}

Supervised classifiers fine-tuned on toxicity labels achieve high accuracy for explicit violations but require task-specific fine-tuning. Prior work has explored geometric structure in embeddings for syntactic parsing [Hewitt \& Manning, 2019], factual knowledge \citep{Petronietal2019}, and bias detection \citep{Bolukbasietal2016}. However, no prior work has tested whether alignment failures form geometric structure in pretrained embeddings or systematically validated base-rates of genuine violations in preference datasets.

This work explicitly tests the geometric manifold hypothesis for alignment failures. The negative result (Cohen's d=0.034) demonstrates pretrained encoders miss safety-specific features despite human detection consistency (κ=0.724), revealing a limitation: general-purpose semantic similarity does not capture alignment distinctions.
